\emph{Шлях звільнення} — це практичний і стислий провідник до духовного звільнення, яке іноді називають пробудженням, просвітленням, самореалізацією або просто баченням того, що є абсолютно Істинним. Неможливо знати, що означають такі слова, як \emph{звільнення} чи \emph{просвітлення}, доки ви самі їх не здійсните. Тож марно міркувати про те, чим є просвітлення; насправді робити це — велика перешкода для того, щоб воно сталося. Один із провідних принципів полягає в тому, що поступово розпізнавати те, що не є абсолютно Істинним, — безмежно цінніше, ніж міркувати про те, чим воно є.

Багато людей вважають, що призначення духовного вчення — давати відповіді на великі питання життя, але насправді все навпаки. Головне завдання доброго духовного вчення — не відповідати на ваші питання, а ставити під сумнів ваші відповіді. Адже саме ваші усвідомлювані й неусвідомлювані передумови та ваші вірування змушують вас бачити поділ там, де існує лише єдність і повнота.

Дійсність, на яку вказують ці вчення, не схована, не є таємною чи далекою. Ви не можете її здобути, заслужити чи зрозуміти. Саме цієї миті Дійсність і повнота — перед очима. Насправді єдине, що існує, щоб бачити, чути, нюхати, куштувати, торкатися чи відчувати, — це Дійсність, або Бог, якщо вам більше до вподоби. Абсолютна повнота оточує вас, де б ви не були. Отже, немає причин цим перейматися — хіба що через той факт, що ми, люди, вже давно загнали себе в такий темний клубок сум'яття й безладу, що ми навряд чи взяли б до уваги існування — а ще менше пережили б самі — божественність усередині нас і навколо нас.

\emph{Шлях звільнення} — це заклик до дії; це те, що ви \emph{робите}. Це \emph{роблення, яке цілком вас розробить}. Якщо ви не \emph{практикуватимете} вчення, якщо не вивчатимете й не застосовуватимете його безстрашно, жодної трансформації не відбудеться. \emph{Шлях звільнення} — не система вірувань. Він цілком практичний.

Читати цю книгу як глядач означало б упустити суть. Бути глядачем легко й безпечно; бути діяльним учасником власного пробудження до Істини — не легко й не безпечно. Шлях попереду непередбачуваний, зобов'язання абсолютне, а результати не гарантовані. Ви справді думали, що може бути інакше?

Якщо ви порівнюватимете \emph{Шлях звільнення} з іншими вченнями або тлумачитимете його крізь призму інших учень, ви неминуче припуститеся хибних тлумачень. Нині, за миттєвого доступу до духовних учень, ця проблема особливо поширена. Люди неправильно тлумачать те, що я кажу, бо пропускають це крізь призму інших духовних учень, які послуговуються подібним словником. Тож я раджу вам підходити до цих учень, не фільтруючи їх попередніми розуміннями розуму.

Жодне духовне вчення не є прямим шляхом до просвітлення. Насправді шляху до просвітлення немає просто тому, що просвітлення завжди присутнє всюди й у кожну мить. Те, що ви \emph{можете зробити}, — це усунути всі ілюзії, особливо ті, які ви найбільше цінуєте й у яких почуваєтеся в безпеці, — ті, що затуляють ваше сприйняття Дійсності. Відмовтеся від своїх ілюзій і свого опору, і Дійсність раптом стане видимою.

\emph{Шлях звільнення} — це ліки для зцілення різних духовних недуг. Так само як ліки самі по собі не є здоров'ям, а засобом його досягти, ці вчення не є Істиною, а засобом її здійснити. Індійський мудрець Рамана Махарші порівнював духовні вчення з колючками, які використовують, щоб вийняти інші колючки. І особисто мені ця аналогія подобається.

Вивчати \emph{Шлях звільнення} — це вивчати самого себе. Вивчати себе не означає додавати до свого мозку більше знань, ідей про себе, а означає усунути всі поняття й характеристики, з якими ви зазвичай ототожнюєте своє «я»: ім'я, расу, стать, заняття, соціальний статус, минуле, а також психологічні судження, які ви виносите про себе. Коли «я» зведено до своєї сутнісної серцевини, усе, що можна сказати, — це \emph{я є; я існую}.

То що ж це за \emph{я}, яке існує?

Це не книга про поліпшення, самовдосконалення чи змінені стани свідомості. Вона про духовне пробудження — перехід від сновидного стану его до неспанного стану поза его, якомога швидше й ефективніше. Подорож не така, як люди її уявляють, і просвітлення не таке, як його зазвичай продають. Я не навчатиму вас досягати екстазу чи безмежного щастя, знаходити свою споріднену душу чи десяти кроків до того, щоб стати мільйонером. Я не вірю ні в оманливу рекламу, ні в хибні обіцянки. Багато духовних шукачів уже живуть на постійній дієті неякісних учень — тих гарних банальностей, приємних на слух, що спричиняють малу або й жодну трансформацію, окрім полегшення незадоволення, притаманного сновидному стану. Якщо вам подобаються такі речі, ця книга не для вас.

Я залишив підказки для здійснення Істини по всій цій книзі, від першої до останньої сторінки. Не припускайте, що найважливіші елементи цих учень легко помітити чи що вони чітко наголошені. Їх вшито в цю книгу так, як нитки перетворюються на тканину, — їх легко не побачити, якщо у вас немає очей, щоб їх бачити, або щирості, щоб їх зрозуміти. Не те щоб я хотів бути неясним --- я роблю все, що можу, щоб \emph{не} бути неясним, --- але Істина не є чимось, що можна глибоко й цілком зрозуміти, якщо подати її розжованою. Такі істини схожі на \emph{фастфуд}: легко доступні, але не задовільні в довгій перспективі.

У нашому сучасному суспільстві ми очікуємо, що все подадуть нам у легко засвоюваних порціях, бажано дуже швидко, щоб ми могли повернутися до свого квапливого життя. Але Істина не погодиться з нашими шаленіми способами уникати Дійсності чи з нашим бажанням отримати ціле за найменших затрат часу й енергії.

Ви візьмете з учень \emph{Шляху звільнення} рівно те, що в них вкладете. Ці вчення потрібно вивчати, обдумувати й утілювати в практику, а не просто читати заради розваги. Один мудрець колись сказав: «Випробування бажання — у звичці відповіді».

Також потрібно розуміти, що \emph{Шлях звільнення} не є ні формою психотерапії, ні розв'язанням усіх викликів, з якими люди стикаються у своїх буднях. Хоча терапевтичні застосування можуть бути потрібними й корисними для деяких людей, вони не є осереддям цих учень.

Пробудження не є ні ліками від усього, що вас гнітить, ні втечею від життєвих труднощів. Таке магічне мислення йде в протилежному напрямку від розгортання Дійсності й є великою перешкодою для того, щоб вона виявилася в найзрілішій формі. Мета цих учень — прокинутися до абсолютної природи Дійсності, а потім утілити її й жити нею максимально повно. Таке пробудження зрештою може принести відчуття глибокого спокою, любові, добробуту, але це лише побічні продукти неспанного стану, а не мета.

До просвітлення веде не пошук дедалі вищих станів щастя й блаженства, а палке прагнення Дійсності та щире незадоволення життям, яке не є цілком справжнім.

\section*{Прокинься або загинь}

Проблеми світу — це загалом людські проблеми, неминучий наслідок егоїстичного сновидіння наяву. Якщо придивитися уважно, усі ознаки наявні й підказують, що ми не лише сновиди наяву, а й сягнули меж божевілля. Певною мірою ми втратили (або принаймні забули) свій дух, і ми наполегливо намагаємося ігнорувати цей факт, бо не хочемо бачити, наскільки ми заціпенілі, наскільки спустошеним насправді є наш стан. Тож ми йдемо навпомацки, рухомі силами, яких не впізнаємо й не розуміємо, а то й зовсім не помічаємо.

Безсумнівно, ми перебуваємо в дуже критичний момент. Наш світ тримається на вельми хисткій рівновазі. Пробудження до Дійсності вже не просто можливість, а необхідність. Ми пливли кораблем омани так далеко, як могли. Ми кинулися в море й тепер опинилися викинутими на спустошений острів. Наші варіанти вичерпалися. «Прокинься або загинь» — це духовний заклик нашого часу. Невже нам потрібно буде більше спонукань, ніж це?

І все ж усе вічно добре — і краще, ніж можна уявити.
